\chapter{Building and Running on the K1}\label{r02:ch:build-run}

This chapter is the practical recipe: how the workspace reaches the board, how waLBerla
is configured and built there, how the benchmark is run, and how to check that the
binary really contains RVV code. All commands are taken from the scripts in
\file{scripts/} and the board logs.

\begin{pitfall}[The board is shared]
  The Banana~Pi is used by others. Everything lives under \file{/home/bananapi/saleh/}:
  the synced workspace (\file{my-e4-ws/}), the Python venv (\file{venv/}), build
  directories (\file{walberla-build*/}), caches (\file{.cache/}) and logs. Long jobs run
  in detached \code{screen} sessions with output to a log file, so they survive a lost
  connection.
\end{pitfall}

\section{Getting the sources to the board}\label{r02:sec:push}

The workspace is edited on a laptop and mirrored to the board with three scripts in the
workspace root:
\begin{description}[style=nextline]
  \item[\code{bash push.to.sh bananapi}] \code{rsync -az --delete} of the workspace to
    \file{/home/bananapi/saleh/my-e4-ws/}, without \file{.git}, the local venv and
    caches. Incremental, about two seconds per push. The remote directory itself is
    kept, so build directories that point into it survive.
  \item[\code{bash run.on.bananapi.sh <command>}] push, then run a command in the remote
    workspace root with the venv's Python first on \code{PATH} and the \ps{} fork on
    \code{PYTHONPATH}. \code{NO_PUSH=1} skips the push.
  \item[\code{bash test.on.bananapi.sh <pytest args>}] push, then run \ps{}' pytest
    natively on the board (with \code{-o addopts=''}, because \ps{}' default options
    require pytest-cov).
\end{description}
Never run two of them concurrently: a push in the middle of a build changes the
sources under it.

\section{One-time setup}\label{r02:sec:setup}

\begin{shcode}
# Python venv on the board (Python 3.12), outside the synced tree
python3 -m venv /home/bananapi/saleh/venv
/home/bananapi/saleh/venv/bin/pip install numpy sympy pytest     # ...and the usual deps
/home/bananapi/saleh/venv/bin/pip install --no-deps pystencilssfg==0.1.1 lbmpy==2.1.1

# Environment for every configure/build (the scripts export these)
export PATH=/home/bananapi/saleh/venv/bin:$PATH
export PYTHONPATH=/home/bananapi/saleh/my-e4-ws/submodules/pystencils/src:\
/home/bananapi/saleh/my-e4-ws/submodules/walberla/sweepgen/src
export XDG_CACHE_HOME=/home/bananapi/saleh/.cache
\end{shcode}
System packages needed: GCC and/or Clang, CMake $\geq$ 3.26 (the board has 3.28), an
MPI implementation (Open MPI 4.1.6), and \code{llvm-objdump} for inspecting binaries.

\section{Configure}\label{r02:sec:configure}

\rneedspace{12\baselineskip}
\begin{shcode}[nolol=false,caption={[Configuring a benchmark build directory]Configuring a benchmark build directory (from \texttt{scripts/walberla\_bench\_all.sh}).},label=r02:lst:configure]
cmake -S /home/bananapi/saleh/my-e4-ws/submodules/walberla \
      -B /home/bananapi/saleh/walberla-build-gcc-f64 \
      -DCMAKE_BUILD_TYPE=Release \
      -DCMAKE_C_COMPILER=gcc -DCMAKE_CXX_COMPILER=g++ \
      -DWALBERLA_DOUBLE_ACCURACY=ON \
      -DWALBERLA_OPTIMIZE_FOR_LOCALHOST=ON \
      -DWALBERLA_ENABLE_V8CORE=ON -DWALBERLA_ENABLE_SWEEPGEN=ON \
      -DWALBERLA_SWEEPGEN_MANAGED_VENV=OFF \
      -DPython_EXECUTABLE=/home/bananapi/saleh/venv/bin/python \
      -DWALBERLA_BUILD_WITH_MPI=ON \
      -DWALBERLA_BUILD_BENCHMARKS=ON -DWALBERLA_BUILD_TESTS=OFF \
      -DWALBERLA_BUILD_TOOLS=OFF -DWALBERLA_BUILD_TUTORIALS=OFF \
      -DWALBERLA_BUILD_EXAMPLES=OFF -DWALBERLA_BUILD_SHOWCASES=OFF -DWALBERLA_BUILD_DOC=OFF
\end{shcode}
For Clang use \code{clang}/\code{clang++}; for single precision
\code{-DWALBERLA_DOUBLE_ACCURACY=OFF}. Each combination needs its own build directory.
The configure output should contain \emph{``Found pystencils Source File Generator
(Version 0.1.1)''}. To build waLBerla's own tests as well, set
\code{-DWALBERLA_BUILD_TESTS=ON}.

\section{Build}\label{r02:sec:build}

\begin{shcode}
cmake --build /home/bananapi/saleh/walberla-build-gcc-f64 --target SweepgenLbmCPU_explicit_cse -j 2
\end{shcode}
Use at most two parallel jobs: the heaviest waLBerla translation units need about
1.3\,GB of memory each, and the board has under 4\,GB. The first target in a fresh build
directory takes about 40 minutes, because it builds the waLBerla core libraries; each
further benchmark variant takes about 5 minutes. The generated sources land in
\file{<build>/apps/benchmarks/SweepgenLbmCPU/_gen/<target>/gen/}.

\section{Run}\label{r02:sec:run}

\begin{shcode}
cd /home/bananapi/saleh/walberla-build-gcc-f64/apps/benchmarks/SweepgenLbmCPU
taskset -c 2 ./SweepgenLbmCPU_explicit_cse 64 200 10        # one core: 64^3 cells, 200 steps
mpirun -np 4 --bind-to core --map-by core ./SweepgenLbmCPU_explicit_cse 64 100 10   # one cluster
mpirun -np 8 --bind-to core --map-by core ./SweepgenLbmCPU_explicit_cse 64 100 10   # whole chip
\end{shcode}
With $N$ processes the domain has $N$ blocks of $64^3$ cells. The K1's eight cores form
two clusters of four, so 4 processes fill one cluster and 8 the chip.

\section{Unattended runs}\label{r02:sec:scripts}

The full measurement matrix takes several hours and runs from scripts started in
\code{screen}:
\begin{center}\small
\begin{tabularx}{\linewidth}{@{}>{\raggedright\arraybackslash}p{5.2cm}>{\raggedright\arraybackslash}X@{}}
  \toprule
  \textbf{Script} & \textbf{What it does}\\
  \midrule
  \file{scripts/walberla_lbm_bench.sh} & in one build directory: build the six variants,
    count RVV instructions in each generated sweep object, run each three times on core
    2, print a summary table\\
  \file{scripts/walberla_bench_all.sh} & configure and run the \{GCC, Clang\} $\times$
    \{f64, f32\} matrix, one build directory each; then waLBerla's
    \code{TestLbmScenarios_CXX}; then a combined summary\\
  \file{scripts/walberla_bench_mpi.sh} & every executable of every configuration with 4
    and 8 MPI processes\\
  \file{scripts/walberla_bench_followup.sh} & rerun failed configurations, then the
    multi-core runs\\
  \file{scripts/walberla_bench_summary.py} & merge logs into one markdown table with
    speedups and checksum agreement\\
  \bottomrule
\end{tabularx}
\end{center}
\begin{shcode}
cd /home/bananapi/saleh && screen -dmS walberla-all bash -c \
    "bash my-e4-ws/scripts/walberla_bench_all.sh > walberla-all.log 2>&1"
\end{shcode}
Each script waits for the previous \code{screen} session to finish before it starts, so
two benchmarks never compete for the cores.

\section{Checking the generated code}\label{r02:sec:check-code}

Three quick checks show whether a build really uses the RVV backend:
\begin{shcode}
G=apps/benchmarks/SweepgenLbmCPU/_gen/SweepgenLbmCPU_explicit_cse/gen
grep -c '__riscv_' $G/SweepgenLbmCPU.cpp                 # intrinsics in the source
grep 'rvv.hpp'     $G/SweepgenLbmCPU.cpp                 # runtime header included
obj=$(find apps/benchmarks/SweepgenLbmCPU/CMakeFiles/SweepgenLbmCPU_explicit_cse.dir \
           -name 'SweepgenLbmCPU.cpp.o' -path '*_gen*')
llvm-objdump -d --no-show-raw-insn $obj | awk '$2 ~ /^v/ {n++} END {print n}'   # vector insns
\end{shcode}
Use \code{llvm-objdump}: GNU objdump 2.42 on the board prints Clang's RVV instructions as
raw \code{.insn} words and would count zero. For the double-precision GCC build the
counts are 665 vector instructions (\code{explicit}), 608 (\code{explicit_cse}) and 138
(\code{autovec}).

\section{Running the pystencils tests}\label{r02:sec:run-tests}

\begin{shcode}
# laptop (x86): codegen tests, cross-compile syntax checks
cd submodules/pystencils
PYTHONPATH=src ../../.venv-rvv/bin/python -m pytest -o addopts='' tests/nbackend/test_rvv_*.py
# board: everything RVV, natively
bash test.on.bananapi.sh -q tests/nbackend/test_rvv_target.py tests/nbackend/test_rvv_codegen.py \
                            tests/nbackend/test_rvv_execution.py tests/nbackend/test_vectorization.py
NO_PUSH=1 bash test.on.bananapi.sh -q -k RISCV_RVV tests/kernelcreation tests/nbackend
\end{shcode}
Every board test JIT-compiles natively, about ten seconds each; filter with \code{-k} or
\code{-x}.

\section{Known nuisances}\label{r02:sec:nuisances}

\begin{itemize}
  \item \textbf{Clock skew.} The board clock is about 46 hours behind the laptop.
        Synced files carry future modification times, so \code{make} warns about clock
        skew and rebuilds more than needed. This is harmless (it never causes stale
        objects), and the clock of a shared machine is left alone.
  \item \textbf{Caches.} lbmpy and \ps{} write caches under \code{XDG_CACHE_HOME};
        keep it inside the user directory.
  \item \textbf{Package versions.} Do not let pip upgrade pystencils-sfg past 0.1.1, and
        never install \ps{} into the venv; it must come from \code{PYTHONPATH}.
\end{itemize}
